% paper/sections/11b1_cls_conservation.tex
%
% §11.3.1  CLS 法の保存性評価：動的非一様格子上での保存型再マッピング
%          （旧 D.2; \section と \label は 11b_component_verification.tex の
%            \subsubsection で提供する）
%      実験スクリプト: experiments/ls_cls_conservation.py
%      結果データ:    results/ls_cls_conservation/

\paragraph{主張．}

本実験が示す主命題は以下の一文である：

\begin{center}
\textbf{%
  界面追従格子を定期的にリフレッシュする環境では，%
  CLS の保存型再マッピングが LS の非保存補間に比べ%
  質量誤差を 1--2 桁低減する．}
\end{center}

LS（符号距離関数 $\phi$）は格子再生成時に区分線形補間を行うため，
補間ごとに質量誤差が蓄積する．
CLS（平滑化 Heaviside $\psi$）は補間後に質量スケーリング補正を施すため，
各リフレッシュで誤差がリセットされる．

\paragraph{支配方程式．}

計算座標 $\xi = i/N$ と物理座標 $x$ の写像を Jacobian $J_i = \mathrm{d}x/\mathrm{d}\xi$ で定義する．
均一流 $u = \mathrm{const}$ の下での各手法の移流方程式：
\begin{equation}
  \text{LS（非保存型）:}\quad
  \frac{\partial \phi}{\partial t} + \frac{u}{J}\,\frac{\partial \phi}{\partial \xi} = 0,
  \qquad
  \text{CLS（保存型）:}\quad
  \frac{\partial \psi}{\partial t} + \frac{1}{J}\,\frac{\partial (u\psi)}{\partial \xi} = 0.
\end{equation}
質量は計算座標上の中点則で評価する：
\begin{equation}
  M = \frac{1}{N}\sum_{i=0}^{N-1} q_i\, J_i
  \qquad \left(= \int_0^1 q(x)\,\mathrm{d}x\right).
\end{equation}

\paragraph{格子・パラメータ．}

界面位置 $x_c(t)$ 近傍に $r$ 倍の細格子を配置する代数的写像：
\begin{equation}
  J(\xi) \propto \frac{1}{1 + (r-1)\exp\!\left(-\frac{(\xi - x_c)^2}{2\sigma^2}\right)},
  \qquad
  \text{正規化：}\frac{1}{N}\sum J_i = 1.
\end{equation}
パラメータ：$N=128$，$r=2$，$\sigma=0.06$，$h_{\min} \approx 1/(rN)$，
$\varepsilon = 3\,h_{\min}$，CFL $= 0.4$．

$K_{\mathrm{refresh}}$ 移流ステップごとに格子を $x_c(t)$ 中心に再生成し，補間を行う：
\begin{itemize}
  \item \textbf{LS}：区分線形補間（質量非保存）
  \item \textbf{CLS}：区分線形補間 $+$ 質量スケーリング補正
    $\psi_i \leftarrow \psi_i^{\mathrm{interp}} \cdot (M_{\mathrm{old}}/M_{\mathrm{new}})$
\end{itemize}

\paragraph{空間離散化と時間積分．}

空間微分には計算座標 $\xi$ 上の Dissipative CCD (DCCD) 演算子を適用する
（付録~\ref{app:weno5}，$\varepsilon_d = 0.05$）．
時間積分は TVD-RK3（Shu--Osher）．
LS のみ $N_{\mathrm{reinit}} = 20$ ステップごとに Godunov 型再初期化を適用する；
CLS は保存型移流により $\tanh$ 形状が維持されるため再初期化を行わない．

\paragraph{結果．}

\paragraph{プロファイル比較．}
$T=1$（1 周期後）のプロファイルを
両手法とも平滑化 Heaviside $\tilde{H}_\varepsilon$ スケールに変換して比較した．
LS は 20 ステップごとの Godunov 型再初期化により SDF 形状を保つため，
$\tilde{H}_\varepsilon(\phi)$ は exact 解とほぼ重なる（$L_2 \approx 4\times10^{-3}$）．
一方 CLS は再初期化を行わず質量保存を優先するため，
遷移幅に $\Ord{10^{-2}}$ 程度のずれが生じる（$L_2 \approx 6\times10^{-2}$）．

\paragraph{質量誤差の時間発展．}
$K=10$ において，LS はリフレッシュのたびに補間起源の質量誤差が蓄積し増大傾向を示す．
CLS はリフレッシュで誤差がリセットされるため振幅が抑制される．

\paragraph{リフレッシュ周期依存性．}
$T=1$ 後の最終相対質量誤差 $|\Delta M/M_0|$ の $K_{\mathrm{refresh}}$ 依存性が
本実験の主要結果である（表~\ref{tab:ls_cls_mass_err}）．
頻繁なリフレッシュ（小 $K$）では CLS が LS に比べ 1--2 桁小さい誤差を示す
（$K=10$: CLS $8.9\times10^{-7}$ vs LS $7.6\times10^{-5}$，85 倍の改善）．
大 $K$ では DCCD 移流誤差が支配的になり CLS の保存性優位が消える
（$K=50$: CLS $8.2\times10^{-4}$ vs LS $7.3\times10^{-4}$，両者同程度）．

表~\ref{tab:ls_cls_mass_err} に全 $K$ での最終質量誤差をまとめる．

\begin{table}[htbp]
  \centering
  \caption{%
    $T=1$ 後の最終相対質量誤差 $|\Delta M/M_0|$（$N=128$，$r=2$，CFL $=0.4$）．}
  \label{tab:ls_cls_mass_err}
  \begin{tabular}{lrrrr}
    \toprule
    手法 & $K=5$ & $K=10$ & $K=20$ & $K=50$ \\
    \midrule
    LS (SDF)     & $8.6\times10^{-5}$ & $7.6\times10^{-5}$ & $1.2\times10^{-4}$ & $7.3\times10^{-4}$ \\
    CLS ($\psi$) & $7.7\times10^{-6}$ & $8.9\times10^{-7}$ & $7.4\times10^{-5}$ & $8.2\times10^{-4}$ \\
    \bottomrule
  \end{tabular}
\end{table}

\paragraph{結論．}

動的非一様格子において CLS の保存型再マッピングはリフレッシュ起源の質量誤差を
1--2 桁低減する．
この優位性はリフレッシュ周期が短いほど顕著であり，
界面追従格子が頻繁に再生成される実用的な計算条件（$K \lesssim 20$）で特に有効である．
リフレッシュ間隔が長い（$K \gg 1$）場合は DCCD 移流起源の非保存誤差が支配的となるため，
適切な $K$ の選択と移流精度の両立が重要である．
